\documentclass[11pt]{article}
\usepackage[a4paper,margin=25mm]{geometry}
\usepackage{amsmath,amssymb,amsthm}
\usepackage[hidelinks]{hyperref}
\newtheorem{proposition}{Proposition}
\newtheorem{corollary}{Corollary}
\title{Aggregate history, hidden lane allocation, and useful common advice\\Conditional analytical results and the remaining obstruction}
\author{Freeway bottleneck review continuation}
\date{17 September 2026}
\begin{document}
\maketitle
\noindent\textbf{Status.} Methodology only. These are original conditional derivations,
not simulations, an adoption of a traffic law, or a proof about native SUMO.
The approved target is an admissibly prepared pair with equal strong aggregate
history and different useful shared desired-speed-advice decisions. We obtain
an information-side result in a free-flow model class and a decision-side
negative result in a fixed-service queue. Neither completes that target.

\section{Information-side construction without deleting history}
Consider three upstream lanes, one of which ends at a fixed downstream lane drop.
Before any traffic reaches a restrictive boundary, let $k_\ell(t,x)$ denote
continuous density in mass/metre and let all lanes have the same free-flow speed
$v_0>0$ in metres/second. Let $r_{\ell j}\geq0$ be conservative lateral transfer
rates in mass/(metre second), with transfers paired on the same longitudinal
coordinate. The conditional free-flow equations are
\begin{equation}
\partial_t k_\ell+\partial_x(v_0k_\ell)
=\sum_{j\ne\ell}r_{j\ell}-\sum_{j\ne\ell}r_{\ell j}.
\label{eq:lane}
\end{equation}
We require admissible, nonnegative states, adequate receiving/storage, and an
identical homogeneous free-flow branch throughout the prefix. Transfers need
not be zero, symmetric, or independent of density. We do not clip or repair a
published transfer law to enforce these assumptions.

Let $K=\sum_\ell k_\ell$, and provide the controller the entire predecision
history of spatial mass and density, mass-weighted mean speed, empty-bin masks,
cumulative aggregate flows across spatial boundaries, source waiting and input
(both intended and actual), together with common geometry, response parameters, advice and
future forcing. Lane-entry allocation is the sole unknown past input in this
construction. There are no connectors occupied during this prefix.

\begin{proposition}[Exact aggregate-prefix invariance]
Starting empty with identical total entry rate $e(t)$ and identical aggregate
boundary conditions, two admissible solutions of \eqref{eq:lane} have identical
declared aggregate histories, regardless of their hidden lane-entry allocations.
This holds on any spatial/time sampling of these observables.
\end{proposition}
\begin{proof}
Pairwise lateral terms cancel, giving
\begin{equation}
\partial_tK+v_0\partial_xK=0,\qquad
K(0,x)=0,\qquad v_0K(t,0)=e(t).
\label{eq:aggregate}
\end{equation}
The transport solution is determined by the same total input and is therefore
identical. Every spatial-bin mass is its integral, every aggregate crossing
flow is $v_0K$, and mean speed in every occupied bin is $v_0$. Empty bins have
the same mask and the same declared empty-speed convention. Choose entry below
the receiving limits of either allocation, so actual and intended entry agree
and source waiting is zero in both. The other known inputs are identical.
No current-state-only reduction or shorter history has been used.
\end{proof}

\subsection*{A reachable hidden contrast in the conditional class}
The result alone does not guarantee that a lane contrast survives. A short-prefix
example can establish that separately. Suppose, additionally, total outgoing
lateral transfer from any lane is bounded by $\lambda$ times its current mass,
where $\lambda<\infty$ is in inverse seconds. This is a hypothesis about a
valid transfer law, not permission to cap an invalid one. Inject the same low
total mass $m$ over $[0,T]$: preparation $a$ through the ending lane and $b$
through a continuing lane. Let no traffic reach the lane-drop boundary or leave
the upstream region during this short prefix. Total mass satisfies $M(t)\leq m$.
For ending-lane upstream mass $E_a,E_b$, the transfer bound gives
\begin{align}
E_a(T)&\geq m-\lambda\int_0^T M(t)\,dt\geq m-\lambda mT,\\
E_b(T)&\leq \lambda\int_0^T M(t)\,dt\leq\lambda mT,\\
E_a(T)-E_b(T)&\geq m(1-2\lambda T)>0\quad\text{if }T<(2\lambda)^{-1}.
\label{eq:contrast}
\end{align}
For $\lambda=0$, any sufficiently short prefix works. Small enough entry rate
keeps the states on the declared free-flow branch. Thus known-empty initialization
and full aggregate history do not, by themselves, rule out a reachable lane
contrast when lane-entry forcing is hidden.

\textbf{Scope of this witness.} The contrast is upstream ending-lane mass at an
early decision, not certified contrast in the last 400 metres of the approved
numerical design. Persistence to that region and postdecision control value
remain unproved. These are continuous-mass observation analogues, not bit-identical
native vehicle-count/trajectory arrays. If lane inputs are also known, this
construction no longer establishes unknown-input ambiguity. If actual free-flow
speed depends on lane/gaps or receiving limits activate, Proposition 1's
premises must be checked again; they cannot be assumed from homogeneous vehicle types.

\subsection*{Traffic-source grounding and non-transfer}
FB0030 gives paired lane conservation and triangular demand in its equations
1--2 and 12--15. In an admissible, demand-active homogeneous free-flow branch,
the paired modifications in equations 20--22 also cancel upon summation over
lanes. This motivates the model class above, but does not certify the complete
published discrete scheme. Its calibrated lanes have different fundamental
diagrams; empty-density ratios, transfer budgets and receiving/capacity branches
remain unresolved. We have not transplanted the continuous short-prefix witness
into that scheme, disabled mandatory/courtesy incentives, or proven its lane
contrast persists. See \href{https://doi.org/10.1080/21680566.2019.1700846}{FB0030}.

\section{Why a hidden contrast is not yet useful advice information}
We now test the tempting completion: attach a fixed bottleneck capacity to each
hidden preparation, and claim this alone makes slowing advice differently useful.
It does not in the following explicitly restricted class.

For each hidden state $s$, model the bottleneck as a work-conserving fluid point
queue, initially empty, with nonnegative locally integrable service rate $c_s(t)$.
The rate may differ between preparations, but is fixed independently of the
subsequent advice. Let $B_{s,u}(t)$ be cumulative arrivals to that queue under
common desired-speed advice trajectory $u$. Let $S_s(t)=\int_0^t c_s(z)\,dz$.
Assume some common feasible trajectory $u^+$ gives arrival dominance:
\begin{equation}
B_{s,u^+}(t)\geq B_{s,u}(t)\quad\text{for all }s,u,t.
\label{eq:arrivals}
\end{equation}
For an independent free-flow approach with monotone desired-speed response,
holding the highest feasible advice supplies such a trajectory. This assumption
is \emph{not} inferred for native following, lane changing or spillback.

Queue departures satisfy the continuous min-plus expression
\begin{equation}
D_{s,u}(t)=\inf_{0\leq r\leq t}
\left\{B_{s,u}(r)+S_s(t)-S_s(r)\right\},\qquad B_{s,u}(0)=0.
\label{eq:queue}
\end{equation}
Use continuous cumulative arrivals here; atomic time-zero batches require an
explicit initial-queue convention and are not needed for this derivation.
Let $I(t)$ be cumulative intended departures for a common finite cohort of mass
$m$, including traffic still waiting at source or on the approach. Define the
complete all-cohort time-spent score through queue exit by
\begin{equation}
J_s(u)=\frac{1}{m}\int_0^\infty\left[I(t)-D_{s,u}(t)\right]dt.
\label{eq:cost}
\end{equation}
Assume every cohort drains and scores are finite; a common fixed downstream
travel time adds an action-independent constant. Any already accrued prefix
cost is likewise fixed within each preparation.

\begin{proposition}[A common best action in the fixed-service class]
Under these assumptions, $J_s(u^+)\leq J_s(u)$ for every state and feasible advice.
Consequently no pair in this class requires disjoint optimal or nonnegative-
tolerance acceptable action sets over the full feasible advice space.
\end{proposition}
\begin{proof}
For the same $S_s$, arrival dominance makes every expression inside the infimum
in \eqref{eq:queue} no smaller for $u^+$. Thus
$D_{s,u^+}(t)\geq D_{s,u}(t)$ pointwise. The intended-input process in
\eqref{eq:cost} is unchanged, so integration gives the score inequality. This
argument counts delayed arrivals on the approach/source; moving waiting outside
the bottleneck cannot create an apparent improvement. It holds separately for
every $s$, including different fixed capacities. All states share $u^+$.
\end{proof}

For a feasible set $\mathcal U$ containing $u^+$ and acceptable tolerance
$\varepsilon_s\geq0$, this means
\begin{align}
\mathcal A_s&=\{u\in\mathcal U:J_s(u)-\inf_{v\in\mathcal U}J_s(v)
\leq\varepsilon_s\},\qquad u^+\in\bigcap_s\mathcal A_s,\\
\inf_{u\in\mathcal U}\max_s
\left[J_s(u)-\inf_{v\in\mathcal U}J_s(v)\right]&=0.
\end{align}
This is a conditional obstruction, not proof that real common speed advice
cannot help. A mechanism that violates arrival dominance or makes realized
service depend on advice/state evolution can escape it. Capacity drop is one
possible mechanism, not a necessary condition for every traffic-control benefit.

\section{The missing mechanism and the action-menu qualification}
The remaining target needs a jointly justified chain: admissible lane/history
collision near the consequential region; shared desired-speed response; endogenous
merge/receiving dynamics under that response; and disjoint useful-action sets
after all slowing/source/recovery costs. Positive lane effects under another
actuator do not complete this chain.

FB0001 provides an aggregate speed-limit/capacity-drop analytical example, but its
equations 5/9 use a discontinuous discharge switch and equation 6 regulates
upstream inflow rather than the approved desired-response zone. FB0038 optimizes
lateral fractions with natural lane changing turned off; it is not shared advice.
FB0039 supplies field plausibility for common speed limits affecting lane
distribution, not a dynamic voluntary-advice closure or paired-history certificate.
See \href{https://arxiv.org/abs/1310.2658}{FB0001},
\href{https://doi.org/10.1177/0361198120948055}{FB0038}, and
\href{https://doi.org/10.1016/j.trc.2017.01.024}{FB0039}.

The prospective nine-action menu has baseline advice 90 km/h and early/late
60-second pulses of $\pm5,\pm10$, reverting to 90. It does not contain holding
100 km/h throughout the decision horizon. Proposition 2 is therefore a result
for a feasible advice space containing the dominating action, not a negative
result for that pulse menu. Restricted-menu distinctions could still occur
without unavoidable continuous-space information loss. Conversely, including
held 100 in a native reference would not assume it is actually best there.
An added dominance stress probe requires explicit menu/budget approval, not
silent expansion. With the existing maximum paired/context/time-step counts,
one added action changes action invocations from 216 to 240 and the total cap
from 480 to 504; preparation counts and other obligations remain unchanged.

\section{Certificate disposition}
\begin{itemize}
\item Established conditionally: exact full aggregate-prefix equality and a
reachable early lane contrast in a homogeneous conservative free-flow class.
\item Established conditionally: fixed, advice-independent work-conserving
service plus arrival dominance cannot produce an unavoidable useful-action split.
\item Not established: the near-drop preparation, native full-history collision,
an admissible shared-advice/merge service closure, or a useful-action split.
\item Not established: minimum sufficient retained state, any gradient correctness,
independent gradient usefulness, field relevance, or generic aggregate sufficiency.
\end{itemize}
We stop at the evidence-route decision. No unsupported lane penalty, shortened
history, categorical moving-mass process, smoothing or saltation exception has
been adopted to bridge the missing link.
\end{document}
